\section{Conclusion}

\project{} treats semantic equivalence as the first-class obstacle to proving
a legacy game's execution.  The initial replay runner demonstrates that the
reconstructed engine can consume real v1.02h witnesses deterministically and
complete full runs, while strict parsing removes several ambient-memory and
format-trust hazards.  The result is not yet a proof system and is not yet an
exact model of the shipped executable.

An experimental shared binary protocol now emits real per-frame subsystem and
root digests with raw floating-point bits, stable identities, and fail-closed
decoding.  Complete repeated Linux runs are byte-identical, and the first broad
ANM experiment already exposed the distinction between allocated memory and
future-live semantic state.  The next concrete milestone is to resolve the
known reuse and dynamic-job omissions, add selected field-level snapshots, and
emit the same schema from the shipped or an exact-reference executable.  This
should turn ``the replay desynchronized'' into a reproducible first-frame,
first-subsystem, first-field diagnosis.

That evidence will guide the projection, but it cannot prove a slice sound.
The intended end state additionally requires a machine-checked commuting
transition and noninterference argument, exact or refined arithmetic semantics,
and a reviewed binding from those definitions to both executables.  Only then
can zkVM performance results support the intended claim.
